\section{Sistemes d'equacions lineals}

Un \textbf{sistema d'equacions lineals} està format per equacions de primer grau en cadascuna de les incògnites, i totes les igualtats s'han de verificar alhora. En forma matricial $A\cdot X = B$:

\begin{equation*}
\begin{pmatrix} a_{11} & a_{12} & \cdots & a_{1n} \\ a_{21} & a_{22} & \cdots & a_{2n} \\ \vdots & \vdots & \ddots & \vdots \\ a_{m1} & a_{m2} & \cdots & a_{mn} \end{pmatrix}\begin{pmatrix} x_1 \\ x_2 \\ \vdots \\ x_n \end{pmatrix} = \begin{pmatrix} b_1 \\ b_2 \\ \vdots \\ b_m \end{pmatrix}
\end{equation*}

La \textbf{matriu ampliada} és $(A \mid B)$, que inclou els termes independents com a columna afegida.

\subsection{Tipus de sistemes}

\begin{itemize}
\item \textbf{Compatible determinat (S.C.D):} té una única solució.
\item \textbf{Compatible indeterminat (S.C.I):} té infinites solucions.
\item \textbf{Incompatible (S.I):} no té solució.
\end{itemize}

Dos sistemes són \textbf{equivalents} si tenen exactament les mateixes solucions.

\section{El mètode de Gauss}

El \textbf{mètode de Gauss} consisteix a transformar la matriu ampliada en una forma escalonada mitjançant operacions elementals de files (intercanviar files, multiplicar per un escalar, sumar un múltiple d'una fila a una altra), fins obtenir zeros a l'esquerra:

\begin{enumerate}
\item Escriure la matriu ampliada $(A \mid B)$.
\item Col\textperiodcentered{}locar a la primera fila l'equació amb coeficient $1$ a la primera incògnita, si és possible.
\item Eliminar la primera incògnita de les files restants.
\item Repetir el procés per a les files i incògnites següents fins obtenir l'escaló de zeros.
\end{enumerate}

\subsection{Exemple S.C.I: sistema compatible indeterminat}

Considerem el sistema:

\begin{equation*}
\begin{cases} x + 2y - z = 1 \\ 2x + 4y - 2z = 2 \\ 3x + 6y - 3z = 3 \end{cases}
\end{equation*}

Matriu ampliada i escalonament:

\begin{equation*}
\left(\begin{array}{ccc|c} 1 & 2 & -1 & 1 \\ 2 & 4 & -2 & 2 \\ 3 & 6 & -3 & 3 \end{array}\right) \xrightarrow{F_2 - 2F_1,\; F_3 - 3F_1} \left(\begin{array}{ccc|c} 1 & 2 & -1 & 1 \\ 0 & 0 & 0 & 0 \\ 0 & 0 & 0 & 0 \end{array}\right)
\end{equation*}

$\operatorname{Rg}(A) = \operatorname{Rg}(A^*) = 1 < 3$ incògnites $\Rightarrow$ S.C.I. Prenent $y = \lambda$, $z = \mu$ (paràmetres lliures):

\begin{equation*}
x = 1 - 2\lambda + \mu, \quad y = \lambda, \quad z = \mu \qquad \lambda, \mu \in \mathbb{R}
\end{equation*}

\section{Teorema de Rouché-Frobenius}

Donat un sistema $m\times n$ amb matriu de coeficients $A$ i matriu ampliada $A^* = (A \mid B)$:

\begin{itemize}
\item $\operatorname{Rg}(A) = \operatorname{Rg}(A^*) = n \Rightarrow$ \textbf{S.C.D} (solució única).
\item $\operatorname{Rg}(A) = \operatorname{Rg}(A^*) < n \Rightarrow$ \textbf{S.C.I} (infinites solucions; $n - \operatorname{Rg}(A)$ paràmetres lliures).
\item $\operatorname{Rg}(A) \neq \operatorname{Rg}(A^*) \Rightarrow$ \textbf{S.I} (sense solució).
\end{itemize}

\section{Sistemes homogenis}

Un sistema és \textbf{homogeni} si tots els termes independents són zero ($B = \vec 0$). En aquest cas $\operatorname{Rg}(A) = \operatorname{Rg}(A^*)$ sempre, de manera que un sistema homogeni és \textbf{sempre compatible}.

\begin{itemize}
\item $\operatorname{Rg}(A) = n \Rightarrow$ S.C.D: l'única solució és la solució trivial $x_1 = x_2 = \cdots = 0$.
\item $\operatorname{Rg}(A) < n \Rightarrow$ S.C.I: infinites solucions (la trivial sempre n'és una).
\end{itemize}

\section{Regla de Cramer}

Per a sistemes S.C.D de $n$ equacions i $n$ incògnites ($|A| \neq 0$), cada incògnita es calcula substituint la columna corresponent de $A$ per la columna $B$ de termes independents:

\begin{equation*}
x_i = \frac{|A_i|}{|A|}
\end{equation*}

Per a un sistema $3\times 3$ amb incògnites $x, y, z$:

\begin{equation*}
x = \frac{\begin{vmatrix} b_1 & a_{12} & a_{13} \\ b_2 & a_{22} & a_{23} \\ b_3 & a_{32} & a_{33} \end{vmatrix}}{|A|},\quad y = \frac{\begin{vmatrix} a_{11} & b_1 & a_{13} \\ a_{21} & b_2 & a_{23} \\ a_{31} & b_3 & a_{33} \end{vmatrix}}{|A|},\quad z = \frac{\begin{vmatrix} a_{11} & a_{12} & b_1 \\ a_{21} & a_{22} & b_2 \\ a_{31} & a_{32} & b_3 \end{vmatrix}}{|A|}
\end{equation*}

\section{Resolució per matriu inversa}

En un sistema S.C.D ($|A| \neq 0$, existeix $A^{-1}$), el sistema $A\cdot X = B$ es pot resoldre aïllant $X$:

\begin{equation*}
A\cdot X = B \;\Longrightarrow\; A^{-1}\cdot A\cdot X = A^{-1}\cdot B \;\Longrightarrow\; X = A^{-1}\cdot B
\end{equation*}

\subsection{Exemple}

Resolem el sistema:

\begin{equation*}
\begin{cases} x + y = 3 \\ 2x - y = 0 \end{cases} \;\Rightarrow\; A = \begin{pmatrix} 1 & 1 \\ 2 & -1 \end{pmatrix},\quad B = \begin{pmatrix} 3 \\ 0 \end{pmatrix}
\end{equation*}

$|A| = (1)(-1) - (1)(2) = -3 \neq 0$, per tant existeix $A^{-1}$:

\begin{equation*}
A^{-1} = \frac{1}{-3}\begin{pmatrix} -1 & -1 \\ -2 & 1 \end{pmatrix} = \begin{pmatrix} \tfrac13 & \tfrac13 \\ \tfrac23 & -\tfrac13 \end{pmatrix}
\end{equation*}

\begin{equation*}
X = A^{-1}\cdot B = \begin{pmatrix} \tfrac13 & \tfrac13 \\ \tfrac23 & -\tfrac13 \end{pmatrix}\begin{pmatrix} 3 \\ 0 \end{pmatrix} = \begin{pmatrix} 1 \\ 2 \end{pmatrix} \;\Rightarrow\; x = 1,\; y = 2
\end{equation*}

\section{Discussió i resolució de sistemes amb paràmetre}

Quan un sistema conté un paràmetre $k$:

\begin{enumerate}
\item Calcular $|A|$ en funció de $k$ i trobar els valors que l'anul\textperiodcentered{}len ($k = a, b, \ldots$).
\item Si $k \neq a, b, \ldots$: $|A| \neq 0 \Rightarrow \operatorname{Rg}(A) = n \Rightarrow$ S.C.D. Resoldre per Cramer o inversa.
\item Per cada valor singular $k = a, b, \ldots$: substituir a $A^*$, escalonar (Gauss) i aplicar el Teorema de Rouché-Frobenius per classificar i resoldre.
\end{enumerate}

\section{Complement: problemes i interpretació geomètrica}

\subsection{Com plantejar un problema amb un sistema}

\begin{enumerate}
\item \textbf{Triar les incògnites} i dir explícitament què representa cadascuna ($x$, $y$, $z$ amb unitats).
\item Traduir cada frase de l'enunciat a una \textbf{equació}. Hi sol haver una equació de totals, una de valors (preus, quantitats\ldots{}) i una relació entre incògnites.
\item Resoldre (Gauss, Cramer o matriu inversa) i \textbf{discutir} si la solució té sentit (ha de ser positiva i, sovint, entera).
\item Respondre en paraules i comprovar a l'enunciat.
\end{enumerate}

\emph{Exemple:} Un cinema ven entrades d'adult (8 \euro{}), joves (5 \euro{}) i infantils (3 \euro{}). Un dia es van vendre $200$ entrades i es van recaptar $1100$ \euro{}. Es sap que hi va haver el doble d'adults que d'infantils. Quantes n'hi va haver de cada tipus?

Sigui $x$ = adults, $y$ = joves i $z$ = infantils.

\begin{equation*}
\begin{cases}x+y+z=200\\8x+5y+3z=1100\\x=2z\end{cases}
\end{equation*}

Substituint $x=2z$: $y=200-3z$ i $16z+5(200-3z)+3z=1100\Rightarrow 4z=100\Rightarrow z=25$. Llavors $x=50$ i $y=125$.

\emph{Comprovació:} $50+125+25=200$ i $400+625+75=1100$ $\checkmark$. \textbf{Hi va haver $50$ adults, $125$ joves i $25$ infantils.}

\subsection{Interpretació geomètrica dels sistemes $3\times3$}

Cada equació d'un sistema de tres incògnites és un \textbf{pla} de l'espai. Resoldre el sistema és trobar la intersecció dels tres plans.

\begin{table}[!htbp]
\centering
\renewcommand{\arraystretch}{1.3}
\begin{tabularx}{\textwidth}{>{\raggedright\arraybackslash}X>{\raggedright\arraybackslash}X}
\toprule
Tipus de sistema & Posició dels tres plans \\
\midrule
S.C.D. (una solució) & Es tallen en \textbf{un punt} \\
S.C.I. (infinites solucions, 1 paràmetre) & Es tallen en una \textbf{recta} (o dos coincidents i un que els talla) \\
S.C.I. (infinites solucions, 2 paràmetres) & Els tres plans són \textbf{coincidents} \\
S.I. (sense solució) & Sense punt comú: tres plans paral\textperiodcentered{}lels, dos paral\textperiodcentered{}lels, o es tallen dos a dos formant un ``prisma'' \\
\bottomrule
\end{tabularx}
\end{table}

\emph{Criteri:} $\operatorname{Rg}(A)$ i $\operatorname{Rg}(A^*)$ determinen el cas (vegeu el tema 10 per a l'estudi detallat de la posició relativa de tres plans).

\section{Complement: més exemples de sistemes resolts}

\emph{Exemple 1 --- Sistema compatible determinat:} Resolem $\begin{cases}x+y+z=6\\2x-y+z=3\\x+2y-z=2\end{cases}$ pel mètode de Gauss. Sumant la primera i la tercera equació, $2x+3y=8$; sumant la segona i la tercera, $3x+y=5$. De la segona, $y=5-3x$ i, substituint, $2x+15-9x=8\Rightarrow x=1$, $y=2$, $z=3$. Amb Rouché--Frobenius: $\operatorname{Rg}A=\operatorname{Rg}A^{*}=3=n$, \textbf{S.C.D.}

\emph{Exemple 2 --- Regla de Cramer sobre el mateix sistema:} El determinant de la matriu de coeficients és $|A|=\begin{vmatrix}1&1&1\\2&-1&1\\1&2&-1\end{vmatrix}=7\neq0$. Aleshores

\begin{equation*}
x=\frac{\begin{vmatrix}6&1&1\\3&-1&1\\2&2&-1\end{vmatrix}}{7}=\frac77=1,\qquad y=\frac{\begin{vmatrix}1&6&1\\2&3&1\\1&2&-1\end{vmatrix}}{7}=\frac{14}{7}=2,\qquad z=\frac{\begin{vmatrix}1&1&6\\2&-1&3\\1&2&2\end{vmatrix}}{7}=\frac{21}{7}=3.
\end{equation*}

\emph{Exemple 3 --- Sistema incompatible:} $\begin{cases}x+y+z=1\\x-y+2z=0\\2x+3z=3\end{cases}$. La suma de les dues primeres equacions és $2x+3z=1$, que contradiu la tercera ($2x+3z=3$). Aleshores $\operatorname{Rg}A=2<\operatorname{Rg}A^{*}=3$: \textbf{S.I.} Geomètricament, els tres plans són secants dos a dos formant un prisma.

\emph{Exemple 4 --- Sistema compatible indeterminat:} $\begin{cases}x+y+z=2\\x-y=0\\2x+z=2\end{cases}$. Les dues últimes donen $y=x$ i $z=2-2x$, i la primera queda $x+x+2-2x=2$, que es compleix sempre. Per tant $\operatorname{Rg}A=\operatorname{Rg}A^{*}=2<3$ (\textbf{S.C.I.}, un paràmetre): $(x,y,z)=(\lambda,\lambda,2-2\lambda)$, que és una recta (feix de tres plans).

\emph{Exemple 5 --- Discussió amb paràmetre:} Discutim $\begin{cases}x+y+z=1\\x+ay+z=a\\x+y+az=a^{2}\end{cases}$ segons $a$. El determinant és $|A|=\begin{vmatrix}1&1&1\\1&a&1\\1&1&a\end{vmatrix}=(a-1)^{2}$. Si $a\neq1$, $\operatorname{Rg}A=\operatorname{Rg}A^{*}=3$: \textbf{S.C.D.} Si $a=1$, les tres equacions són $x+y+z=1$ ($\operatorname{Rg}A=\operatorname{Rg}A^{*}=1$): \textbf{S.C.I.} amb dos paràmetres (un pla).

\section{Complement: el mètode de Gauss pas a pas}

El mètode de Gauss transforma el sistema en un d'equivalent \textbf{escalonat} amb tres tipus d'operacions: intercanviar equacions, multiplicar una equació per un nombre $\neq0$ i sumar a una equació un múltiple d'una altra. S'escriu només la \textbf{matriu ampliada} $A^{*}=(A\,|\,b)$ i es fan zeros per sota de cada pivot.

\emph{Exemple:} Resolem $\begin{cases}x+2y+z=8\\2x-y+z=3\\3x+y-2z=-1\end{cases}$.

\begin{equation*}
\left(\begin{array}{ccc|c}1&2&1&8\\2&-1&1&3\\3&1&-2&-1\end{array}\right)\xrightarrow[F_3\to F_3-3F_1]{F_2\to F_2-2F_1}\left(\begin{array}{ccc|c}1&2&1&8\\0&-5&-1&-13\\0&-5&-5&-25\end{array}\right)\xrightarrow{F_3\to F_3-F_2}\left(\begin{array}{ccc|c}1&2&1&8\\0&-5&-1&-13\\0&0&-4&-12\end{array}\right).
\end{equation*}

De baix cap amunt: $-4z=-12\Rightarrow z=3$; $-5y-3=-13\Rightarrow y=2$; $x+4+3=8\Rightarrow x=1$. La solució és $(1,2,3)$ i el sistema és S.C.D. ($|A|=20\neq0$).

\textbf{Més equacions que incògnites.} Si hi ha $4$ equacions i $3$ incògnites, s'escalona igual; les files de zeros indiquen equacions redundants i una fila $(0\ 0\ 0\,|\,c)$ amb $c\neq0$ indica S.I.

\emph{Exemple:} $\begin{cases}x+y+z=6\\x-y+2z=5\\2x+y-z=1\\x+z=4\end{cases}$: s'obté $\operatorname{Rg}A=\operatorname{Rg}A^{*}=3=n$, de manera que és S.C.D. amb solució $(1,2,3)$ i una equació sobra.

\section{Complement: discussió completa d'un sistema amb paràmetre}

Discutim i resolem segons $m$: $\begin{cases}x+y+mz=1\\x+my+z=m\\mx+y+z=m^{2}\end{cases}$.

\textbf{Pas 1 --- determinant.} Per Sarrus, $|A|=\begin{vmatrix}1&1&m\\1&m&1\\m&1&1\end{vmatrix}=-(m-1)^{2}(m+2)$, que s'anul\textperiodcentered{}la per $m=1$ i $m=-2$.

\textbf{Pas 2 --- cas general.} Si $m\neq1$ i $m\neq-2$, $\operatorname{Rg}A=\operatorname{Rg}A^{*}=3=n$: \textbf{S.C.D.} Per exemple, per a $m=0$ la solució és $\bigl(\tfrac12,\tfrac12,-\tfrac12\bigr)$.

\textbf{Pas 3 --- cas $m=1$.} Les tres equacions són $x+y+z=1$. $\operatorname{Rg}A=\operatorname{Rg}A^{*}=1<3$: \textbf{S.C.I.} amb dos paràmetres, $(x,y,z)=(1-\lambda-\mu,\lambda,\mu)$. Els tres plans són coincidents.

\textbf{Pas 4 --- cas $m=-2$.} El sistema és $x+y-2z=1$, $x-2y+z=-2$, $-2x+y+z=4$. Sumant les tres equacions s'obté $0=3$: $\operatorname{Rg}A=2<\operatorname{Rg}A^{*}=3$, \textbf{S.I.} Els plans es tallen dos a dos formant un prisma.

\emph{Resum:}

\begin{table}[!htbp]
\centering
\renewcommand{\arraystretch}{1.3}
\begin{tabular}{lccc}
\toprule
Valor de $m$ & $\operatorname{Rg}A$ & $\operatorname{Rg}A^{*}$ & Tipus \\
\midrule
$m\neq1,\,-2$ & $3$ & $3$ & S.C.D. \\
$m=1$ & $1$ & $1$ & S.C.I. (2 paràmetres) \\
$m=-2$ & $2$ & $3$ & S.I. \\
\bottomrule
\end{tabular}
\caption{Discussió del sistema segons $m$}
\end{table}

\section{Complement: sistemes homogenis amb paràmetre}

Un sistema homogeni $AX=0$ és \textbf{sempre compatible} (té la solució trivial). Si $|A|\neq0$ només té la solució trivial (S.C.D.); si $|A|=0$ té infinites solucions (S.C.I.).

\emph{Exemple:} Per a quins valors de $k$ té solucions no trivials $\begin{cases}x+y+kz=0\\x+ky+z=0\\kx+y+z=0\end{cases}$? El determinant és $-(k-1)^{2}(k+2)$.

\begin{itemize}
\item Si $k\neq1,-2$: només la solució trivial $(0,0,0)$.
\item Si $k=1$: queda $x+y+z=0$, un pla; $\operatorname{Rg}=1$ i la solució té dos paràmetres.
\item Si $k=-2$: $\operatorname{Rg}=2$ i la solució té un paràmetre: $(x,y,z)=\lambda(1,1,1)$.
\end{itemize}

\section{Complement: problemes de sistemes d'equacions}

\textbf{Problema 1 (edats).} Una mare, la seva filla i la seva néta tenen entre les tres $60$ anys. La mare té el doble d'anys que la filla i la filla el triple que la néta. Quines edats tenen?

Siguin $x$, $y$, $z$ les edats de la mare, la filla i la néta: $x+y+z=60$, $x=2y$, $y=3z$. Substituint, $6z+3z+z=60\Rightarrow z=6$. Per tant, \textbf{la néta té $6$ anys, la filla $18$ i la mare $36$}.

\textbf{Problema 2 (mescles).} Es volen mesclar tres tipus de cafè de $6$, $8$ i $10$ \euro{}/kg per obtenir $100$ kg de mescla a $8{,}2$ \euro{}/kg. Si es volen $10$ kg més del primer que del segon, quants kg cal de cadascun?

Siguin $x$, $y$, $z$ els kg de cada tipus: $x+y+z=100$; $6x+8y+10z=820$; $x=y+10$. Substituint: $2y+z=90$ i $14y+10z=760$. De la primera, $z=90-2y$ i $14y+900-20y=760\Rightarrow y=\tfrac{70}{3}$. La solució no és entera, però sí vàlida: $x=\tfrac{100}{3}$, $y=\tfrac{70}{3}$, $z=\tfrac{130}{3}$ kg (positives).

\emph{Comprovació:} $\tfrac{100+70+130}{3}=100$ i $\tfrac{600+560+1300}{3}=820$ $\checkmark$.

\section{Complement: quin mètode triar}

\begin{table}[!htbp]
\centering
\renewcommand{\arraystretch}{1.3}
\begin{tabularx}{\textwidth}{>{\raggedright\arraybackslash}X>{\raggedright\arraybackslash}X}
\toprule
Situació & Mètode recomanat \\
\midrule
Sistema $3\times3$ numèric & Gauss \\
Cal només saber si és S.C.D. & $|A|\neq0$ \\
Sistema amb paràmetre & $|A|$, valors crítics, Rouché--Frobenius \\
Sistema $n\times n$ que es resol diverses vegades amb termes diferents & Matriu inversa $X=A^{-1}b$ \\
Només una incògnita & Cramer \\
Més equacions que incògnites & Gauss (Rouché--Frobenius) \\
Homogeni & $|A|=0$ per a solucions no trivials \\
\bottomrule
\end{tabularx}
\caption{Quin mètode triar per resoldre un sistema}
\end{table}

\section{Complement: relació amb la geometria}

Cada equació lineal amb tres incògnites és un \textbf{pla} de l'espai; resoldre un sistema $3\times3$ és estudiar la intersecció de tres plans. Els casos S.C.D., S.C.I. i S.I. corresponen a les vuit posicions relatives de tres plans (vegeu la figura de la posició relativa de tres plans del tema 10): tres plans que es tallen en un punt (S.C.D.), en una recta o coincidents (S.C.I.), i els casos sense punts comuns (S.I.).
